\documentclass[a4paper,10.5pt]{article}
\usepackage[utf8]{inputenc}
\usepackage[T1]{fontenc}
\usepackage[french]{babel}
\usepackage[left=1cm,right=1cm,top=.5cm,bottom=0cm]{geometry}
\usepackage{enumitem}
\usepackage{hyperref}
\usepackage{xcolor}
\usepackage{titlesec}
\usepackage{fontawesome5}
\usepackage{lmodern}
\usepackage[normalem]{ulem}

% --- Couleurs Radiall ---
\definecolor{primary}{RGB}{0, 48, 135}
\definecolor{secondary}{RGB}{80, 80, 80}

% --- Configuration des liens ---
\hypersetup{
    colorlinks=true,
    linkcolor=primary,
    filecolor=primary,
    urlcolor=primary,
}

% --- Formatage des titres ---
\titleformat{\section}{\large\bfseries\color{primary}\uppercase}{}{0em}{}[\titlerule]
\titlespacing{\section}{0pt}{8pt}{5pt}

% --- Commandes personnalisées ---
\newcommand{\resumeItem}[1]{\item\small{#1}}

\begin{document}

% --- EN-TÊTE ---
\begin{center}
    {\Large \textbf{Loïc Aron MBASSI EWOLO}} \\ \vspace{4pt}
    \textbf{\large Ingénieur Développement Produits \& Tests Embarqués} \\
    \textit{\small Disponible dès \textbf{Septembre 2026} en alternance (12 mois)} \\ \vspace{4pt}
    \small
    \faMapMarker\ Cergy, France (Mobile IDF) \quad
    \faPhone\ (+33) 7 59 52 33 98 \quad
    \faEnvelope\ \href{mailto:loicaron.mbassiewolo@ensea.fr}{loicaron.mbassiewolo@ensea.fr} \\ \vspace{2pt}
    \faLinkedin\ \href{https://www.linkedin.com/in/loic-mbassi/}{linkedin.com/in/loic-mbassi} \quad
    \faGithub\ \href{https://github.com/Nameless0l}{github.com/Nameless0l} \quad
    \faGlobe\ \href{https://mbassiloic.tech}{mbassiloic.tech}
\end{center}

\vspace{-6pt}

% --- PROFIL ---
\noindent\small{Prototypage hardware (soudure, capteurs IMU/Flex), automatisation de tests en \textbf{Python} et développement embarqué en \textbf{C/C++} sur STM32 et Nvidia Jetson. Double diplôme \textbf{ENSEA} / \textbf{Polytechnique Yaoundé} en systèmes embarqués et électronique. Je souhaite contribuer au développement et à la validation des modules de transmission sans contact chez \textbf{Radiall}.}

% --- FORMATION ---
\section{Formation}
\begin{itemize}[leftmargin=*, label=\tiny$\bullet$, nosep]
    \item \textbf{ENSEA (École Nationale Supérieure de l'Électronique et de ses Applications)} \hfill \textit{Cergy | 2025 -- 2027} \\
    \small{Double diplôme d'Ingénieur : \textbf{Systèmes Embarqués}, Traitement du Signal, IA.} \\
    \textit{\small{Cours : Conception numérique FPGA/VHDL, Protocoles de communication, Électronique analogique/numérique.}}
    \vspace{2pt}
    \item \textbf{École Polytechnique de Yaoundé (1\textsuperscript{ère} école d'ingénieur CEMAC)} \hfill \textit{Yaoundé | 2021 -- 2025} \\
    \small{Diplôme d'Ingénieur de Conception (Master 2) : \textbf{Mathématiques Appliquées}, Génie Logiciel, Systèmes.}
\end{itemize}

% --- COMPÉTENCES TECHNIQUES ---
\section{Compétences Techniques}
\begin{itemize}[leftmargin=*, nosep]
    \item \small{\textbf{Embarqué \& Électronique :} C/C++, STM32, Arduino, Raspberry Pi, Soudure, Capteurs (IMU, Flex), Protocoles (I2C, SPI, UART).}
    \item \small{\textbf{Automatisation \& Scripts :} \textbf{Python}, Bash, Automatisation de tests, Traitement de données, PyInstaller.}
    \item \small{\textbf{IHM \& Développement :} Qt/Tkinter, React, API REST, Git, Docker, Linux.}
    \item \small{\textbf{Outils \& Conception :} MATLAB/Simulink, VHDL (FPGA), Onshape (CAO 3D), KiCad (PCB).}
    \item \small{\textbf{Langues :} Français (Natif), Anglais (Professionnel/Technique).}
\end{itemize}

% --- PROJETS & EXPÉRIENCES ---
\section{Projets \& Expériences}
\begin{itemize}[leftmargin=*, label={-}]

\item \textbf{Harmony Gloves -- Gant Instrumenté pour Langue des Signes} \hfill \textit{2024} \\
\textit{Laboratoire IOTA, Polytechnique Yaoundé} \hfill \textit{STM32, Arduino, C, Capteurs IMU/Flex}
\vspace{-2pt}
    \begin{itemize}[leftmargin=0.4cm, nosep, label=\tiny$\bullet$]
        \item \small{Participation au prototypage d'un gant connecté : \textbf{soudure} des composants, intégration capteurs IMU et flex sur PCB.}
        \item \small{Développement firmware en \textbf{C} pour l'acquisition et la fusion de données capteurs en temps réel.}
        \item \small{Campagne de \textbf{prise de données} expérimentales et validation du prototype par tests fonctionnels.}
    \end{itemize}
\vspace{3pt}

\item \textbf{Upgrade Jetson -- Robotique Autonome (Challenge Défense CoHoMa)} \hfill \textit{2025 -- En cours} \\
\textit{ENSEA, collaboration armée française} \hfill \textit{C/C++, Docker, YOLOv10, Lidar 3D, SLAM}
\vspace{-2pt}
    \begin{itemize}[leftmargin=0.4cm, nosep, label=\tiny$\bullet$]
        \item \small{Intégration et traitement de données \textbf{Lidar 3D} (VLP16) et caméra de profondeur sur \textbf{Nvidia Jetson}.}
        \item \small{Développement embarqué en \textbf{C/C++} avec contraintes temps réel et optimisation mémoire.}
        \item \small{Déploiement reproductible via \textbf{Docker} et modélisation 3D sur Onshape.}
    \end{itemize}
\vspace{3pt}

\item \textbf{Fault Detection \& Fault-Tolerant Control (Drone)} \hfill \textit{2025} \\
\textit{Projet académique ENSEA} \hfill \textit{MATLAB/Simulink, Contrôle automatique}
\vspace{-2pt}
    \begin{itemize}[leftmargin=0.4cm, nosep, label=\tiny$\bullet$]
        \item \small{Modélisation dynamique d'un drone et conception d'un \textbf{banc de test} de détection de défauts (capteurs/actionneurs).}
        \item \small{Simulation de scénarios de pannes réalistes et \textbf{analyse comparative} des performances sous Simulink.}
    \end{itemize}
\vspace{3pt}

\item \textbf{Outil d'Analyse Vidéo Automatisé (YOLOv8)} \hfill \textit{2025} \\
\textit{Projet personnel} \hfill \textit{Python, YOLOv8, OpenCV, PyInstaller}
\vspace{-2pt}
    \begin{itemize}[leftmargin=0.4cm, nosep, label=\tiny$\bullet$]
        \item \small{Développement d'un \textbf{script Python} de détection automatique et extraction de séquences vidéo.}
        \item \small{Optimisation mémoire pour traitement longue durée. Compilation en \textbf{exécutable (.exe)} pour distribution.}
    \end{itemize}
\vspace{3pt}

\item \textbf{Stage Recherche IoT -- TinyML sur Edge Computing} \hfill \textit{Oct. 2023 -- Jan. 2024} \\
\textit{IoT Lab, Polytechnique Yaoundé} \hfill \textit{Python, Arduino, Raspberry Pi, TF Lite, C}
\vspace{-2pt}
    \begin{itemize}[leftmargin=0.4cm, nosep, label=\tiny$\bullet$]
        \item \small{Optimisation de modèles ML pour \textbf{hardware limité} (quantization, pruning) avec contraintes temps réel.}
        \item \small{Déploiement sur \textbf{Arduino/Raspberry Pi} et validation par mesures sur banc de test.}
    \end{itemize}
\vspace{3pt}

\end{itemize}

% --- DISTINCTIONS ---
\section{Distinctions \& Leadership}
\begin{itemize}[leftmargin=*, nosep]
    \item \small{\textbf{1\textsuperscript{ère} Place} -- IndabaX Africa 2025 (IA Santé) | \textbf{1\textsuperscript{ère} Place} -- CITS National Hackathon (Urban Waste, 75 équipes).}
    \item \small{\textbf{Chef Cellule Projets \& IA} (Polytechnique Yaoundé) : animation workshops techniques, +50\% membres.}
    \item \small{\textbf{Certifications :} Supervised ML (DeepLearning.AI/Stanford), Python for Data Science (IBM).}
\end{itemize}

\end{document}
